% ===================================================================
%  TABELLA nr. 8 — La racolta.
%  Traduction 2026 d'apres texto/io, controlee sur texto/fr, texto/ca
%  et texto/oc. Consigne : voir texto/rm/10-tabella-01.tex.
%
%  LA NOTE (*) DE CETTE PAGE EST LA PLUS UTILE DE TOUT LE LIVRET, ET
%  C'EST GUIGNON QUI L'ECRIT CONTRE LUI-MEME.
%
%  Il vient d'intituler son tableau « La Rekolto » et il s'arrete
%  aussitot : le mot n'est pas precis, dit-il — moisson de lin ?
%  vendange ? cueillette des pommes ? Il propose « mesono », paru
%  dans Mondo XIV p. 242, et note que l'Academie ne l'a pas encore
%  adopte. AUTREMENT DIT, IL LUI MANQUE UN MOT GENERAL, ET IL LE
%  DEMANDE PAR ECRIT A UN COMITE.
%
%  LE ROMANCHE A EXACTEMENT LE MEME MANQUE ET NE LE SENT PAS. Il n'a
%  pas de mot unique pour l'operation abstraite ; il a un mot par
%  recolte — on « sea » le foin, on « meda » le grain, on
%  « vendemia » la vigne — et surtout il a fait de ces travaux SON
%  CALENDRIER : le tableau 3 l'a montre, juin est « zercladur », le
%  mois ou l'on bine, et juillet « fanadur », le mois ou l'on fane.
%
%  D'OU LA LIGNE QUE LES TABLEAUX 3, 7 ET 8 TRACENT ENSEMBLE : LE MOT
%  GENERAL EST UN BESOIN DU LECTEUR, LE MOT PARTICULIER EST UN BESOIN
%  DU TRAVAILLEUR. Le tableau 3 le montrait sur les mois, le 7 sur
%  les betes — sept series derivees en ido contre six racines pleines
%  en romanche —, le 8 le montre sur le travail lui-meme, et cette
%  fois c'est l'auteur qui le dit tout haut. Une langue faite pour
%  etre imprimee a besoin d'un mot qui couvre ; une langue faite pour
%  etre criee d'un champ a l'autre a besoin d'un mot qui distingue.
%
%  ON GARDE DONC « racolta » AU TITRE, avec la note, parce que la
%  note EST le sujet de la page ; mais dans le corps chaque recolte
%  reprend son verbe et son nom.
%
%  ECART DECLARE DU BLOC t08-c2-03-1. Le fac-simile ido y ecrit
%  « gobleto 13) » sans la parenthese ouvrante. La coquille est du
%  compositeur de l'ido et non du texte romanche : la parenthese est
%  retablie, et renvoji.py passe l'ecart sans rien dire. UNE
%  TRADUCTION N'EST PAS UNE TRANSCRIPTION.
% ===================================================================

%%K t08-noto-1 noto
\VUnotes{143.2mm}{%
\textit{(*) Perquai che quel pled n'è betg precis: è quai la racolta dal glin, la
vendemia, la racolta dals mails? Forsa fiss quai bun da duvrar «}
\VUgras{mesono} \textit{», proponì en Mondo XIV, p. 242. Ma enfin ussa n'è quella
nova ragisch betg adoptada da l'Academia ed ella spetga las discussiuns tenor las
reglas usitadas dals idists.}%
}

%%K t08-apar-1 apar
\VUsaut{5.0mm}
\VUtitre{15.0pt}{60}{TABELLA nr. 8}
\VUsaut{3.0mm}
\VUcentre{13.2pt}{90}{\textit{Emprima scena.}}
\VUsaut{3.0mm}
\VUpk{10.2pt}{40}{La racolta (*).}
\VUsaut{4.0mm}
\VUpk{10.2pt}{40}{Ils aspects da la champagna.}
\VUsaut{3.0mm}

%%K t08-c1-01-1 p
1. --- Cun la \VUgras{stad} è l'aspect da la champagna puspè midà. Il sem confidà al
surtg ha germinà; ina verda plantetta è sortida da la terra ed ha fatg spias. La
grina è sa furmada, lura è ella madirada, ed ussa basta d'ir a la retscheiver.

%%K t08-c1-02-1 p
2. --- Las \VUgras{flurs champestras} èn avertas. \VUgras{Flurs da furment}
\textsuperscript{(1)}, \VUgras{papavers} \textsuperscript{(2)} e \VUgras{digitalas}
\textsuperscript{(3)} mischeidan a la nianza uniforma dals chomps da furment il cler
contrast da lur \VUgras{corollas} frestgas.

%%K t08-c1-03-1 p
3. --- Ils arbers da fritga sa èn ornads cun fegls; la fritga è madirada. Il
pitschen \VUgras{ceriser} \textsuperscript{(4)} è plain bellas \VUgras{ceriaschas}
\textsuperscript{(5)}, ch'ils uffants racoglian e maglian cun grond plaschair.

%%K t08-c1-04-1 p
4. --- Uss po il muvel passentar tut il di en l'aria libra.

%%K t08-c1-04-2 p
Ils \VUgras{agnels} \textsuperscript{(6)} èn creschids e daventads bellas \VUgras{nursas}
\textsuperscript{(7)} e bels \VUgras{muntuns} \textsuperscript{(8)} cun in dens vel da
launa. Els pascentan cun quietezza l'\VUgras{erva} da la \VUgras{pastira}
\textsuperscript{(9)} smaltgada da \VUgras{margaritas}.

%%K t08-c1-04-3 p
Prest vegn ins a tundrar \VUgras{quellas} bes-chas per avair lur \VUgras{launa}, da la
quala ins fa il drap da nossas vestgadiras.

%%K t08-tit-2 sub
\VUsaut{5.0mm}
\VUpk{10.2pt}{40}{Il pastur.}
\VUsaut{3.4mm}

%%K t08-c1-05-1 p
5. --- Il \VUgras{pastur} \textsuperscript{(10)} para da betg als dar fitg quità. El sa
fa mo quità da la tempesta che passa a l'orizont, ed il \VUgras{guardian da las
nursas} \textsuperscript{(13)}, che avischinescha ina \VUgras{flascha}
\textsuperscript{(14)} a sias leftgas, para anc pli senza quità.

%%K t08-c1-06-1 p
6. --- Il chaud è opprimant; perquai vegn era il \VUgras{chaun} \textsuperscript{(15)} a
stizzar sia said; el lapa l'aua frestga da l'\VUgras{auatg} \textsuperscript{(16)} e
spaventa ina paupra \VUgras{rana} \textsuperscript{(17)} che sa zuppava en l'erva da la
riva. Quella sigla en l'aua clera, nua che flureschan las \VUgras{ninfeas}
\textsuperscript{(18)}.

%%K t08-c1-07-1 p
7. --- Ina \VUgras{batticua} \textsuperscript{(19)} sa bogna sper la \VUgras{funtanetta}
\textsuperscript{(20)}, che sprizza tranter duas crappas e sa zuppa sut ils
\VUgras{buschs spinus} \textsuperscript{(21)} e las \VUgras{spinablancas}
\textsuperscript{(22)}.

%%K t08-tit-3 sub
\VUsaut{5.0mm}
\VUpk{10.2pt}{40}{La racolta.}
\VUsaut{3.4mm}

%%K t08-c1-08-1 p
8. --- Per ils uffants da champagna è la racolta dal furment ina epoca fitg
divertenta. Els \VUgras{sa rulan}\textsuperscript{(24)} allegramain sin ils
\VUgras{mantuns da palia} \textsuperscript{(23)}, els gidan ils \VUgras{sieaders}
\textsuperscript{(26)} e, entant che quels sean il \VUgras{chomp da furment}
\textsuperscript{(27)} cun lur \VUgras{faultschs} \textsuperscript{(25)} bain agiavgiads
cun il \VUgras{cot} \textsuperscript{(30)}, rimnan els il furment e al ligian en
\VUgras{gerbas} \textsuperscript{(31)} u en \VUgras{faschs} \textsuperscript{(32)}.

%%K t08-c1-09-1 p
9. --- Mintgatant permetta ins als pli gronds da tagliar cun il \VUgras{faultschin}
\textsuperscript{(12)} ils \VUgras{stips d'aur} \textsuperscript{(28)} che portan las
pesantas \VUgras{spias} \textsuperscript{(29)} plainas da grina; ed ils pli pitschens,
survegliond il \VUgras{muvel} \textsuperscript{(33)} che pascenta en la \VUgras{pastira}
\textsuperscript{(34)} sut l'umbriva dal grond \VUgras{tigl} \textsuperscript{(35)},
tschiffan cun lur \VUgras{rait} \textsuperscript{(36)} las bellas \VUgras{farfaglias}.

%%K t08-c1-09-2 p
Intgins rampignan schizunt sin il \VUgras{char} \textsuperscript{(37)} per metter en
urden las gerbas ch'ins als dat si cun ina \VUgras{furtga} \textsuperscript{(38)}.

%%K t08-c1-10-1 p
10. --- Sin tut la \VUgras{plaunira} \textsuperscript{(39)}, nua che sgolan las
\VUgras{lodolas} \textsuperscript{(41)}, regia l'animaziun. Tuts èn occupads cun la
racolta. Ma, entant ch'ils pauprs cuntinueschan a duvrar ils vegls utensils,
\VUgras{faultschs, faultschins} e \VUgras{flagels} \textsuperscript{(40)}, laschan ils
ritgs possessurs sear lur furment u lur \VUgras{sejel} \textsuperscript{(43)} cun ina
\VUgras{maschina da sear} \textsuperscript{(42)}. Lura, senza stuair purtar las gerbas en
il granari, fa ins gronds \VUgras{mantuns da gerbas} \textsuperscript{(44)}.

%%K t08-c1-11-1 p
11. --- Alura maina ins qua ina \VUgras{maschina da bater} \textsuperscript{(47)}, ch'in
\VUgras{locomobil} \textsuperscript{(45)} mova cun ina \VUgras{curriaglia da transmissiun}
\textsuperscript{(46)}, ed uschia vegn la grina separada da la \VUgras{palia}, senza
bandunar il chomp nua ch'ella è vegnida semnada.

%%K t08-tit-4 sub
\VUsaut{5.0mm}
\VUpk{10.2pt}{40}{La tempesta.}
\VUsaut{3.4mm}

%%K t08-c1-12-1 p
12. --- Savens vegnan \VUgras{temporals} \textsuperscript{(53)} en il temp da la racolta.
L'emprim vegn ins ad udir dalunsch ils grivlims dal \VUgras{tun}; in \VUgras{vent ferm da
vest} \textsuperscript{(54)} cumenza a suflar e curva schuffland ils pli gronds arbers
dal \VUgras{guaud} \textsuperscript{(49)}. Nairs \VUgras{nivels} \textsuperscript{(56)}
assaglian il tschiel; els vegnan traversads dals zickzacks sbargliants da la
\VUgras{saetta} \textsuperscript{(55)}. La \VUgras{plievgia} e mintgatant la
\VUgras{granida} cumenzan a crudar, sfrantschond las racoltas prontas per esser
seadas.

%%K t08-c1-13-1 p
13. --- Savens accenda la saetta ina construcziun. Uschia è l'onn passà il tetg dal
\VUgras{malin da vent} \textsuperscript{(50)} vegnì reducì en tschendra e duas da sias
\VUgras{alas} \textsuperscript{(51)} sfrantschadas.

%%K t08-c1-14-1 p
14. --- Suenter intginas uras renascha la quietezza. In splendurus \VUgras{artg dal
tschiel} \textsuperscript{(52)} cumpara a l'orizont e la natira reprenda sia vita
interrutta durant intgins mumaints. Ils purs, ch'èn sa refugiads en las
\VUgras{chamonas} \textsuperscript{(48)}, cumenzan danovamain a lavurar e sa dattan
curaschusamain fadia da reparar ils donns ch'els han gist patì.

%%K t08-tit-5 sub
\VUsaut{6.0mm}
\VUcentre{13.2pt}{90}{\textit{Segunda scena.}}
\VUsaut{3.6mm}

%%K t08-tit-6 sub
\VUpk{10.2pt}{40}{La chatscha. --- La vendemia.}
\VUsaut{1.4mm}
\VUpk{10.2pt}{40}{Il pestgar.}
\VUsaut{4.0mm}

%%K t08-c2-01-1 p
1. --- Enturn la fin d'\VUgras{avust} u la mesa da \VUgras{settember}, tenor las
regiuns, vegn l'avertura da la chatscha. Apaina ch'ils emprims radis dal sulegl han
fatg sortir ils \VUgras{luscherts} \textsuperscript{(2)} da lur fora, sa metta il
\VUgras{chatschader} \textsuperscript{(5)} en viadi cun ses \VUgras{chauns}
\textsuperscript{(4)}.

%%K t08-c2-02-1 p
2. --- El ha mess sia solida \VUgras{vestgadira da chatscha} \textsuperscript{(9)} da
drap gross, el ha mess als peis \VUgras{gambieras} da corom, cun las qualas el vegn a
pudair chaminar en l'erva umida nua che sa zuppan ils \VUgras{botschs}
\textsuperscript{(1)}; el ha prendì ses \VUgras{sclop} \textsuperscript{(6)} e ses
\VUgras{satg da chatscha} \textsuperscript{(10)}, e qua parta el.

%%K t08-c2-03-1 p
3. --- L'ardur da la persecuziun al maina en il mez d'ina vigna nua che la vendemia è
fitg viva. Ils uffants dal vignerun, che survegliaschan uschiglio las chauras sin la
via, las laschan oz pascentar a bel plaschair e, suenter avair ligià ad in pal in
vegl \VUgras{boc} \textsuperscript{(3)} che n'avess betg mancà da profitar da sia
libertad per scappar, sa dattan els era lur part dal plaschair. In tschiffa ina rapa
en in \VUgras{chanaster da vendemia} \textsuperscript{(12)} e sa diverta a laschar
currer il \VUgras{suc} \textsuperscript{(14)} en in magiel \textsuperscript{(13)} per far
most (vin betg fermentà). In auter sa cuvra il chau cun ina \VUgras{curuna da fegls}
\textsuperscript{(15)} ch'el ha gist plitgà.

Sper els vegnan \VUgras{passers} \textsuperscript{(16)} senza vargugna schizunt avant la
trolla per rubar las rapas.

%%K t08-c2-04-1 p
4. --- Entant ch'ils uffants sa divertan, lavuran ils umens. Ils \VUgras{vendemiaders}
\textsuperscript{(18)} portan las rapas en il chaler cun \VUgras{brentas}
\textsuperscript{(17)}; in \VUgras{buttatscher} \textsuperscript{(19)} repara ils
\VUgras{barigls} \textsuperscript{(20)}, verifitgescha sche las \VUgras{duvas}
\textsuperscript{(22)} ed ils \VUgras{funds} \textsuperscript{(21)} èn en urden, e
remplazza ils \VUgras{tschercels} \textsuperscript{(23)} rutts. El ha era lavurà las
\VUgras{tinas} \textsuperscript{(25)}, en las qualas ins versa las \VUgras{rapas}
\textsuperscript{(26)} per las laschar fermentar.

%%K t08-c2-05-1 p
5. --- Cur ch'la fermentaziun è finida, trai ins il vin ed ins metta ils rests sin la
\VUgras{trolla} \textsuperscript{(28)}, e stgatschond adina pli fitg cun la \VUgras{vit}
\textsuperscript{(29)} squitscha ins ora il suc, che curra en ina \VUgras{tinetta}
\textsuperscript{(32)}, ed uschia obtegna ins anc \VUgras{vin} \textsuperscript{(31)}.

%%K t08-tit-8 sub
\VUsaut{6.0mm}
\VUpk{10.2pt}{40}{Biera e sider.}
\VUsaut{3.4mm}

%%K t08-c2-06-1 p
6. --- Vi dal mir dal \VUgras{chaler} \textsuperscript{(27)} rampigna in rom da
\VUgras{luvel} \textsuperscript{(30)}. Qua è el mo ina planta ornamentala, ma en
tscherts pajais, nua ch'la vit è fitg rara, vegn il luvel culvità per far la
\VUgras{biera}.

%%K t08-c2-07-1 p
7. --- Il pitschen \VUgras{mailer} \textsuperscript{(33)} è chargià da mails che vegnan
a dar, cur ch'els vegnan madirs, in excellent \VUgras{sider}. En sia \VUgras{chabgia}
\textsuperscript{(34)} guardan il \VUgras{canarin} \textsuperscript{(35)} e la
\VUgras{cardinera} \textsuperscript{(36)} cun egls envidius ils passers che petulan
libramain.

%%K t08-c2-08-1 p
8. --- Fin a la fin da la vista, sin il pendent da las collinas, èn alineads
\VUgras{pals da vit} \textsuperscript{(40)} che portan las bellissimas \VUgras{chips da
vit} \textsuperscript{(39)}, ornadas cun \VUgras{rapas} pesantas.

%%K t08-c2-08-2 p
Durant tut l'onn ha il \VUgras{vignerun} \textsuperscript{(42)} lavurà per tgirar sia
\VUgras{vigna} \textsuperscript{(38)} ed uss vegn el recumpensà per sia fadia cun ina
bella racolta. Las \VUgras{vendemiadras} \textsuperscript{(41)} emplenan las tinas
messas sin il char tratg da \VUgras{muls} \textsuperscript{(43)}, e tuts èn allegers.

%%K t08-tit-9 sub
\VUsaut{5.0mm}
\VUpk{10.2pt}{40}{Il pestgar.}
\VUsaut{3.4mm}

%%K t08-c2-09-1 p
9. --- Il medem è cun ils \VUgras{pestgaders cun la cauna} \textsuperscript{(44)}, ch'èn
vegnids a s'installar da la damaun a l'ur da l'aua cun lur \VUgras{caunas da pestgar}
\textsuperscript{(45)}.

%%K t08-c2-09-2 p
Il \VUgras{pesch} \textsuperscript{(47)} tschiffa l'\VUgras{am} immediatamain cur ch'els
tuffan lur \VUgras{fil} \textsuperscript{(46)} en l'aua, ed apaina ch'els han gì temp
da renovar l'\VUgras{esca}, sa fa gia ina nova victima al chau dal fil.

%%K t08-c2-09-3 p
Tuttina tschiffa il \VUgras{pestgader cun la rait} \textsuperscript{(48)}, che bitta da
sia \VUgras{barca} \textsuperscript{(49)} la \VUgras{rait da bittar} en il mez dal
\VUgras{flum} \textsuperscript{(50)}, dapli peschs.

%%K t08-c2-10-1 p
10. --- L'atun è la stagiun da las bunas spassegiadas: a pe, \VUgras{a chaval}
\textsuperscript{(52)}, \VUgras{cun il velo} \textsuperscript{(53)}, \VUgras{cun l'auto}
\textsuperscript{(54)}. Las \VUgras{vias} \textsuperscript{(51)} èn nettas ed ins
profitescha dals ultims bels dis, perquai che qua sa radunan gia las
\VUgras{randulinas} \textsuperscript{(56)} sin ils tetgs per sgolar cun tut lur alas
vers pajais pli chauds. Il fegliom dal vegl \VUgras{nuscher} \textsuperscript{(57)}
daventa mellen, las \VUgras{nuschs} crudan, ed ils auts \VUgras{arbers} disponids en
gruppa sco in \VUgras{bugliet} \textsuperscript{(58)} sin l'\VUgras{autezza}
\textsuperscript{(59)} vegnan prest a perder lur fegls.

%%K t08-c2-11-1 p
11. --- Il \VUgras{sulegl} \textsuperscript{(60)} sa auza pli tard, ils dis sa
scursanan, in fraid bastantamain pitschient cumenza ad esser sentì la \VUgras{damaun},
e qua bandunan gia svols da \VUgras{becassas} \textsuperscript{(61)} noss climas pauc
ospitalers.

\VUsaut{6.0mm}
\VUfilet{22mm}
